\scanpage{175}
\mspage{33}

$\unsure{X|\mathscr X}$ est filtrante, $e_{X^\wedge}$ est limite inductive filtrante de foncteurs représentables. On sait d'ailleurs qu'on peut réindexer par \unsure{une|un} catégorie \uline{ordonnée} filtrante (cf SGA4, I, \S\,8.1.6 (c'est dû à P. Deligne ?).

On a prouvé $a\Longrightarrow b$, prouvons $b\Longrightarrow a$. Supposons donc

\[
e_{X^\wedge}=\varinjlim_{i\in I}x_i,
\]

où $I$ est un \inserted{(petit)} ens. ordonné filtrant, et $I\longrightarrow X$ un foncteur. \deleted{Soit \unsure{Prouvons}} Prouvons PS1$'$ \inserted{(forme \unsure{droite|duale})} i-e. pour $a,b$ dans $X$, existe

\diagram{\begin{tikzcd}[ampersand replacement=\&]
a\arrow[r] \& x\\ b\arrow[ur]\&
\end{tikzcd}}

. En effet, \deleted{il ex\unsure{.}} comme \deleted{\unsure{$\varinjlim\operatorname{Hom}(a,x)$}} $\varinjlim_I\operatorname{Hom}(a,x_i)=e$, il s'ensuit que $\operatorname{Hom}(a,x_i)\ne\varnothing$ pour $i$ grand, de $=$ $\operatorname{Hom}(b,x_i)\ne\varnothing$, donc il suffit de prendre

\diagram{\begin{tikzcd}[ampersand replacement=\&]
a\arrow[r] \& x_i\\ b\arrow[ur]\&
\end{tikzcd}}

. Prouvons $\mathrm{PS2}_{\unsure{d|a}}$, i.e. que pour double flèche $a\overset{u}{\underset{v}{\rightrightarrows}}b$, existe $b\unsure{\mathbin{-}|\longrightarrow}x$ qui égalise. Considérons,

\deleted{\unsure{Hom} $i\to j$, $\exists$ \unsure{compatible}, \unsure{$x_i\to x$}}

\diagram{\begin{tikzcd}[ampersand replacement=\&]
\deleted{\xi_i\in\operatorname{Hom}(b,x_i)}\arrow[d]\arrow[r,shift left=.6ex,"{\deleted{u^*}}"]\arrow[r,shift right=.6ex,"{\deleted{v^*}}"']\&\deleted{\operatorname{Hom}(a,x_i)}\arrow[d]\\
\deleted{\lim\operatorname{Hom}(b,x_i)}\arrow[r,"{\deleted{\sim}}"]\&\deleted{\lim\operatorname{Hom}(a,x_i)=e}
\end{tikzcd}}

% archive sheet number: 175
